% Narrow primary-source audit.  This file is editorial only and is not a
% production theorem or proof.

\section*{CLD smoothness versus local logarithmic affineness}

Chambert--Loir--Ducros define an ordinary smooth function on an analytic
space to be, locally for the ordinary topology, a function of the form
\[
  \varphi(\log|f_1|,\ldots,\log|f_r|),
  \qquad f_i\in\mathcal O_X^\times,
  \quad \varphi\in C^\infty(\mathbf R^r);
\]
see \cite[\S0.2.2]{CLD25}.  Consequently, an expression
\[
  c+\sum_i a_i\log|f_i|,\qquad a_i\in\mathbf R,
\]
already proves ordinary CLD smoothness.  Neither rationality of the
coefficients nor algebraization of the analytic units is needed for this
implication.

There is a stronger rigidity statement when the function is also
piecewise linear.  By \cite[Proposition~12.5.4]{CLD25}, if \(u\) is both
paralinear and smooth, then, at every \(x\), there exist an integer
\(n\geq1\), an ordinary analytic neighbourhood \(U\) of \(x\), a real
constant \(c\), and \(g\in\mathcal O_X^\times(U)\) such that
\[
                 n u|_U=c+\log|g|.
\]
If the absolute value of the ground field is nontrivial and the germ
\((X,x)\) is strict, the proposition also allows \(c=0\).  Its proof uses
\cite[Proposition~1.9.16]{CLD25}: on a polyhedral subspace, a function that
is both smooth and piecewise linear is locally the restriction of one
ambient affine form.  The affine form can be chosen with rational slopes;
clearing denominators and multiplying the corresponding analytic units
gives the displayed one-unit formula.  The analytic-space notions of
affine and integral-paralinear functions occur in
\cite[\S12.4.6 and \S12.5.2]{CLD25}.

Thus two assertions must not be conflated:
\begin{enumerate}
\item an ambient real-affine tropical expression pulls back to a real
  affine combination of logarithms of analytic units, which directly
  proves CLD smoothness;
\item a function already known to be paralinear and smooth becomes the
  logarithm of one analytic unit only after multiplication by a possibly
  nontrivial integer, by Proposition~12.5.4.
\end{enumerate}
In particular, the integer \(n\) should not be silently replaced by \(1\).
Conversely, \(c+n^{-1}\log|g|\) is smooth and paralinear, but only its
\(n\)-fold multiple is automatically integral-paralinear in the sense of
\S12.5.2.

\section*{Passage from CLD smooth coordinates to an ordinary GK chart}

Suppose that finitely many scalar coordinates \(u_j\) are CLD-smooth.
Concatenating their finitely many analytic-unit presentations yields one
analytic moment map, and smooth composition shows directly that every
smooth function of the \(u_j\), in particular \(F(u_1,\ldots,u_r)\), is
CLD-smooth.  This step uses only the definition in \S0.2.2.

An additional comparison is required if one asserts a presentation in an
ordinary algebraic tropical chart.  Under the standing hypotheses that
\(K\) is algebraically closed, complete, and nontrivially valued and that
\(X\) is algebraic, \cite[Proposition~7.2]{Gub16} says that an analytic
moment map on an open subset of \(X^{\mathrm{an}}\) can, near any prescribed
point, be replaced by an algebraic moment map to the same split torus with
exactly the same tropicalization.  Applying this once to the product
analytic moment preserves the smooth expression.  Factoring the resulting
algebraic moment through the canonical moment of a very affine open, as in
\cite[\S4.12]{Gub16}, gives the ordinary algebraic tropical chart used by
\cite[Definition~8.1 and Definition~9.9]{GK17}.

It is therefore correct and economical to cite Proposition~7.2 explicitly
at this point.  It is not logically mandatory to cite that proposition by
number if one instead cites an already packaged comparison: for example,
\cite[Remark~5.4]{Gub16} identifies the CLD sheaf of differential forms
with the algebraic-tropical-chart sheaf and explicitly derives the converse
from Proposition~7.2, the analytic-moment comparison lemma of CLD, and
\S4.12.  Likewise, \cite[Proposition~8.4 and \S8.3]{GK17} packages the
analytic-moment-to-algebraic-chart passage, while
\cite[Remark~9.16]{GK17} directly states that smooth metrics are
delta-metrics.  What is not justified is the bare inference from the CLD
definition to ``smooth in the ordinary algebraic-chart sense'' without one
of these comparison citations.
